\subsection{可和方法}
\label{sec:ch1-summability-methods}

为了从 Fourier 系数恢复函数 $f$, 最好能找到不同于取 Fourier 级数部分和极限的方法, 因为正如已经看到的, 后一种方法并不总是有效.

一种这样的办法是 Cesàro 可和法, 即取部分和的算术平均的极限. 众所周知, 若 $\lim a_k$ 存在, 则
\[
 \lim_{k\to\infty}\frac{a_1+\cdots+a_k}{k}
\]
也存在且取相同的值. 定义
\[
\begin{aligned}
 \sigma_Nf(x)
 &=\frac1{N+1}\sum_{k=0}^{N}S_kf(x)\\
 &=\int_0^1f(t)\frac1{N+1}\sum_{k=0}^{N}D_k(x-t)\,dt\\
 &=\int_0^1f(t)F_N(x-t)\,dt,
\end{aligned}
\]
其中 $F_N$ 是 Fejér 核,
\[
 F_N(t)=\frac1{N+1}\sum_{k=0}^{N}D_k(t)
 =\frac1{N+1}\left(\frac{\sin(\pi(N+1)t)}{\sin(\pi t)}\right)^2.
\]
$F_N$ 具有下列性质:
\[
 F_N(t)\ge0,
\]
\hypertarget{chapter-01-page-022}{}
\begin{equation}
\label{eq:ch1-fejer-kernel-properties}
 \|F_N\|_1=\int_0^1F_N(t)\,dt=1,
 \qquad
 \lim_{N\to\infty}\int_{\delta<|t|<1/2}F_N(t)\,dt=0
 \quad\text{if }\delta>0.
\end{equation}
因为 $F_N$ 为正, 它的 $L^1$ 范数等于其积分, 即为 $1$. Dirichlet 核并非如此: 它的积分等于 $1$ 是因为正部与负部相消, 而它的 $L^1$ 范数随 $N$ 趋于无穷.

\begin{Theorem}
\label{thm:ch1-fejer-norm-convergence}
若 $f\in L^p$, $1\le p<\infty$; 或者 $f$ 连续且 $p=\infty$, 则
\[
 \lim_{N\to\infty}\|\sigma_Nf-f\|_p=0.
\]
\end{Theorem}

\begin{Proof}
因为 $\int F_N=1$, 由 Minkowski 不等式,
\[
\begin{aligned}
 \|\sigma_Nf-f\|_p
 &=\int_{-1/2}^{1/2}\|f(\cdot-t)-f(\cdot)\|_pF_N(t)\,dt\\
 &\le\int_{|t|<\delta}\|f(\cdot-t)-f(\cdot)\|_pF_N(t)\,dt
 +2\|f\|_p\int_{\delta<|t|<1/2}F_N(t)\,dt.
\end{aligned}
\]
当 $1\le p<\infty$ 时,
\[
 \lim_{t\to0}\|f(\cdot-t)-f(\cdot)\|_p=0;
\]
当 $p=\infty$ 且 $f$ 连续时, 同一极限也成立. 因而选择适当的 $\delta$, 可以使第一项任意小. 固定 $\delta$ 后, 由 \eqref{eq:ch1-fejer-kernel-properties}, 第二项趋于 $0$.
\end{Proof}

\begin{Corollary}
\label{cor:ch1-trigonometric-density-uniqueness}
\begin{enumerate}[label=(\arabic*)]
\item 当 $1\le p<\infty$ 时, 三角多项式在 $L^p$ 中稠密;
\item 若 $f$ 可积且对每个 $k$ 都有 $\widehat f(k)=0$, 则 $f$ 恒等于零.
\end{enumerate}
\end{Corollary}

第二种可和方法是把 Fourier 级数看成单位圆(复平面中)上形式极限
\begin{equation}
\label{eq:ch1-fourier-series-unit-disk}
 u(z)=\sum_{k=0}^{\infty}\widehat f(k)z^k
 +\sum_{k=-\infty}^{-1}\widehat f(k)\overline z^{|k|},
 \qquad z=re^{2\pi i\theta}.
\end{equation}
由于 $\{\widehat f(k)\}$ 是有界序列, 此函数在 $|z|<1$ 上有定义. 它可改写为
\[
 u(re^{2\pi i\theta})
 =\sum_{k=-\infty}^{\infty}\widehat f(k)r^{|k|}e^{2\pi ik\theta}
 =\int_{-1/2}^{1/2}f(t)P_r(\theta-t)\,dt,
\]
\hypertarget{chapter-01-page-023}{}
其中
\[
 P_r(t)=\sum_{k=-\infty}^{\infty}r^{|k|}e^{2\pi ikt}
 =\frac{1-r^2}{1-2r\cos(2\pi t)+r^2}
\]
是 Poisson 核. Poisson 核具有与 Fejér 核类似的性质:
\[
 P_r(t)\ge0,
\]
\begin{equation}
\label{eq:ch1-poisson-kernel-properties}
 \int_0^1P_r(t)\,dt=1,
 \qquad
 \lim_{r\to1^-}\int_{\delta<|t|<1/2}P_r(t)\,dt=0
 \quad\text{if }\delta>0.
\end{equation}
因此可以证明一个与定理~\ref{thm:ch1-fejer-norm-convergence} 类似的结果.

\begin{Theorem}
\label{thm:ch1-poisson-norm-convergence}
若 $f\in L^p$, $1\le p<\infty$; 或者 $f$ 连续且 $p=\infty$, 则
\[
 \lim_{r\to1^-}\|P_r*f-f\|_p=0.
\]
\end{Theorem}

由于函数 $u$ 在 $|z|<1$ 上调和, 它是 Dirichlet 问题
\[
 \Delta u=0\quad\text{if }|z|<1,
 \qquad
 u=f\quad\text{if }|z|=1
\]
的解, 其中边界条件按定理~\ref{thm:ch1-poisson-norm-convergence} 的意义理解.

在第~\ref{chap:hardy-littlewood-maximal-function}~章, 我们将研究 $\sigma_Nf(x)$ 和 $P_r*f(x)$ 的几乎处处收敛性.

